\documentclass[12pt, a4paper]{article}
\usepackage[utf8]{inputenc}
\usepackage[spanish]{babel}
\usepackage{amsmath}
\usepackage{amsfonts}
\usepackage{amssymb}
\usepackage{graphicx}
\usepackage{geometry}
\usepackage{booktabs}
\usepackage{hyperref}

\usepackage[table]{xcolor}
\usepackage{caption}
\usepackage{array}

\usepackage[numbers]{natbib}

\title{
    {Trabajo PID: Seguimiento 1}\\
    {\large Identificacion de Dinero Uruguayo}
}

\author{Federico Vareika, Guillermo Golpe, Tomas Hazan}
\date{2026}

\begin{document}

\maketitle

\newpage

\section*{Descripcion}
Este trabajo plantea el desarrollo de un sistema robusto de reconocimiento de
billetes de pesos uruguayos basado en el marco teórico de Hasanuzzaman et al. 
\cite{hasanuzzaman2012robust}.
A diferencia de los métodos globales, este enfoque utiliza la extracción de 
características locales en regiones de interés (ROI) específicas, lo que 
garantiza una mayor fiabilidad ante billetes deteriorados, rotos o parcialmente 
visibles. 

Al centrarse en fragmentos distintivos de la moneda nacional, el sistema permite
una identificación precisa de las denominaciones incluso en condiciones de uso 
real adversas.

Para la ejecución técnica, se realizará una comparativa entre los algoritmos 
SURF y ORB. El interés de este contraste radica en sus diferencias fundamentales 
de procesamiento: mientras SURF ofrece una gran estabilidad mediante descriptores 
de punto flotante y distancias euclidianas, ORB destaca por su eficiencia al 
utilizar descriptores binarios y la distancia de Hamming. Esta distinción 
permite evaluar si la velocidad y el bajo consumo de recursos de ORB compensan 
la robustez tradicional de SURF en el contexto de la moneda uruguaya. Se ha 
prescindido del método SIFT ya que al ser técnicamente similar a SURF, su 
inclusión aportaría menos valor analítico que el contraste directo entre un 
método tradicional de alta precisión y la agilidad extrema de un descriptor 
binario como ORB.

Dado que no existe material público de pesos uruguayos con oclusiones y arrugas,
se construirá una base de datos propia para validar el sistema. Como referencia
metodológica, se tomará el repositorio de Carlos Costa \cite{carlosmccosta}; 
sin embargo, dado que dicho código está desarrollado en C++, nuestra
implementación se realizará íntegramente desde cero en Python con OpenCV. Este 
cambio de lenguaje y la construcción manual nos permitirán adaptar la lógica a
las texturas y marcas de seguridad específicas de nuestros billetes, utilizando 
el repositorio original únicamente como guía conceptual para el flujo de trabajo.

\newpage

\section*{Objetivos Especificos}

\begin{itemize}
    \item Implementar el marco teórico de componentes: Adaptar la metodología de 
        Hasanuzzaman et al. (2012) para realizar un reconocimiento basado en regiones 
        de interés (ROI) en lugar de un análisis global del billete.

    \item Realizar una comparativa algorítmica: Evaluar el rendimiento, la precisión y 
        la velocidad de procesamiento entre SURF (basado en descriptores de punto 
        flotante) y ORB (basado en descriptores binarios).

    \item Construir material propio: Crear una base de datos original de moneda uruguaya 
        que contemple variables críticas como oclusiones, arrugas, rotaciones y 
        cambios de iluminación.

    \item Desarrollar una implementación nativa en Python: Traducir y adaptar la lógica 
        conceptual del repositorio de Carlos Costa (originalmente en C) al ecosistema 
        de Python y OpenCV, asegurando un control total sobre el código y su 
        adaptación a los pesos uruguayos.

    \item Garantizar la robustez del sistema: Validar que el sistema sea capaz de 
        identificar correctamente las denominaciones incluso cuando los billetes están 
        deteriorados o parcialmente visibles.

\end{itemize}

\section*{Planificación Inicial}

\begin{itemize}
    \item Decisión del tema (10hs)
    \item Investigación y profundización del tema elegido (20hs)
    \item Desarrollo de los seguimientos del proyecto (25hs)
    \item Preparación de base de datos (imágenes) a utilizar (5 hs)
    \item Desarrollo de las metodologías (30hs)
    \item Documentación y conclusiones de los resultados obtenidos (30hs)
    \item Preparación de la presentación y defensa final (20hs)
\end{itemize}

Total: 150hs, lo que concluye en 50hs por persona

\newpage

\begin{center}
\renewcommand{\arraystretch}{1.5}
{\rowcolors{2}{black!5}{black!10}
    \begin{tabular}{| c | c | c | m{6cm} |} 
    \hline 
    
    \textbf{Fecha} & 
    \textbf{Tiempo (hs)} & 
    \textbf{Miembro} &
    \textbf{Actividad Realizada}\\

    \hline 

    11/3 &
    3 &
    Tomás Hazan &
    Decisión del tema \\

    11/3&
    3&
    Guillermo Golpe&
    Decisión del tema \\

    11/3&
    3&
    Federico Vareika&
    Decisión del tema \\

    12/3&
    4&
    Federico Vareika&
    Investigación de el funcionamiento del método SURF y revisión de la implementación de repositorio existente \\

    14/3&
    2&
    Tomás Hazan&
    Primera aproximación para entendimiento del método ORB \\

    14/3&
    2&
    Guillermo Golpe&
    Primera aproximación para entendimiento del método SIFT  \\

    15/3&
    1&
    Todos los miembros&
    Puesta en común de la información recaudada para metodologías \\

    16/3&
    2&
    Guillermo Golpe&
    Desarrollo de descripción del primer seguimiento \\

    16/3&
    2&
    Tomás Hazan&
    Desarrollo de objetivos y planificación preliminar del primer seguimiento \\

    16/3&
    1&
    Federico Vareika&
    Transcripción de bosquejos a informe Latex para primer seguimiento \\

    \hline 
\end{tabular}}
\end{center}

\newpage 



\bibliographystyle{plain} 
\bibliography{seguimiento1/Referencias}

\end{document}

